%% kapa: exam=1 round=21 year=2010 period=2 form=A subject=appraisal_law subject_ko=부동산관계법규 seq=14 number=54
%% source: Q-Net 감정평가사 기출문제 (cst00309) artlSeq=5024181 file=extracted/1008/5024181/2010년 제21회 감정평가사 1차시험 기출문제/2교시.hwp
%% conversion: kordoc 4.24.51 -> tools/gen1.py (auto); review=proofread-gpt-sol; flags=none
\begin{question}{54}
국토의 계획 및 이용에 관한 법령상 도시계획시설사업 시행자의 토지 등의 수용 및 사용에 관한 설명으로 옳지 않은 것은?
\begin{choices}
  \item 행정청뿐만 아니라 비행정청인 도시계획시설사업 시행자도 도시계획시설사업에 필요한 토지ㆍ건축물을 수용 또는 사용할 수 있다.
  \item 시행자는 사업시행을 위하여 특히 필요하다고 인정되는 때에는 도시계획시설에 인접한 토지를 일시 사용할 수 있다.
  \item 도시계획시설사업에 관한 실시계획의 고시가 있은 때에는「공익사업을 위한 토지 등의 취득 및 보상에 관한 법률」에 따른 사업인정 및 그 고시가 있었던 것으로 본다.
  \item 시행자는 도시계획시설사업에 필요한 토지에 정착된 물건에 관한 소유권 외의 권리를 수용 또는 사용할 수 있다.
  \item 시행자는 토지소유자 등과 협의가 성립하지 않을 때에는 도시계획시설사업에 관한 실시계획의 인가 고시가 있은 날부터 1년 이내에 토지수용위원회에 재결신청을 하여야 한다.
\end{choices}
\end{question}
\answer{5}
